\chapter{Método}

% Guardrails do capitulo (fonte: docs/living-paper/chapter_structure_freeze.md)
% Objetivo: documentar o metodo de forma reproduzivel no pipeline real.
% Escopo obrigatorio: dados, features, TFT, protocolo experimental, baselines,
% metricas (com politica raw vs post-guardrail), explicabilidade e analytics store.
% Fora de escopo: metodo centrado em CNN para classificacao de sentimento.

\section{Desenho da pesquisa}
Esta pesquisa adota abordagem quantitativa, aplicada e de natureza experimental-computacional, com foco em previsão de séries temporais financeiras sob controle de validade metodológica. O desenho experimental foi estruturado para permitir comparação reprodutível entre configurações de features e hiperparâmetros, com avaliação fora da amostra e registro auditável dos artefatos de execução. Em vez de tratar o desenvolvimento como manual de software, as decisões de engenharia foram formuladas como escolhas metodológicas para reduzir vieses de avaliação, aumentar rastreabilidade e preservar a reprodutibilidade dos resultados \cite{PINEAU2021REPRODUCIBILITY}.

O fluxo geral foi organizado em etapas sequenciais: ingestão de dados, engenharia de features, construção de dataset, treinamento, inferência e atualização analítica. A estratégia prioriza validade temporal e comparabilidade estatística, de modo que interpretações finais não dependam de reruns não controlados nem de agregações heterogêneas de coortes.

\section{Dados e pipelines de processamento}
Os dados utilizados combinam informações de mercado (OHLCV), notícias processadas para sinal de sentimento agregado, indicadores técnicos e variáveis fundamentalistas, com recorte temporal definido em configuração declarativa por ativo. A ingestão e o pré-processamento seguem ordem causal, com padronização temporal em UTC para minimizar ambiguidades de calendário e risco de vazamento temporal em joins e agregações.

Do ponto de vista de persistência, o projeto utiliza camadas de persistência imutáveis e auditáveis para armazenamento canônico dos dados de execução e uma camada analítica derivada e reconstruível para consolidação de métricas e artefatos de comparação. Essa separação permite preservar evidências primárias de experimento ao mesmo tempo em que fornece visões analíticas consolidadas para decisão, sem necessidade de reprocessar todo o treinamento a cada nova análise.

\section{Engenharia de features}
A engenharia de atributos foi estruturada por famílias de \textit{features}: (i) \textit{baseline} de preço e volume, (ii) indicadores técnicos, (iii) sentimento agregado e dinâmica de sentimento, e (iv) variáveis fundamentalistas e razões derivadas. A seleção de conjuntos de \textit{features} é versionada e rastreada por assinatura, permitindo comparar experimentos em nível de configuração com controle explícito sobre quais variáveis efetivamente participaram de cada execução.

Para preservar causalidade, o cálculo das \textit{features} obedece janelas estritamente retrospectivas e políticas \textit{as-of} para variáveis de divulgação não diária. Para fundamentos, observa-se a data real de publicação ou disponibilidade da informação, e não apenas a data de fechamento do período contábil; para sentimento, observa-se o \textit{timestamp} de coleta ou publicação e o corte da sessão usada para construir o alvo. Foi adotada também governança de \textit{warmup} por \textit{feature}, com contabilização do período mínimo necessário para validade estatística de cada atributo antes do uso em treino. Essa estratégia reduz risco de imputação indevida no início da série e melhora a consistência entre treino e inferência.

A composição do candidato confirmatório do trabalho segue o princípio de poda mínima baseada em princípio: parte-se de todas as famílias de \textit{features} disponíveis (PRICE, TECHNICAL, SENTIMENT, FUNDAMENTAL) e aplica-se um único conjunto de regras, declarado antes da observação dos resultados, para remoção de variáveis redundantes ou inelegíveis. Em conformidade com \texttt{docs/00\_overview/STRATEGIC\_DIRECTION.md} (§4.2), as regras de poda são: (i) remoção de \textit{features} com correlação par-a-par superior a \(0{,}95\); (ii) remoção de \textit{features} com mais de \(30\%\) de valores ausentes no conjunto de treino; (iii) remoção de derivadas redundantes da mesma quantidade (por exemplo, \texttt{MA\_5} e \texttt{EMA\_5}, mantida apenas uma sob critério \textit{a priori} documentado); e (iv) remoção de qualquer \textit{feature} sem disponibilidade temporal comprovada no momento da inferência. Por princípio metodológico, a lista final de \textit{features} e o conjunto exato de regras aplicadas são congelados no pré-registro versionado antes da rodada confirmatória, eliminando seleção pós-observação como fonte de variação. A escolha do candidato \textit{all-features} sob poda mínima reflete também a evidência exploratória registrada em \texttt{docs/07\_reports/living-paper/evidence\_log.md} (2026-04-24), segundo a qual a escolha do conjunto de \textit{features} não foi o \textit{driver} dominante de desempenho pontual nas rodadas exploratórias: a variância entre configurações dentro do mesmo conjunto foi materialmente maior que a variância entre conjuntos distintos.

\section{Configuração do modelo TFT}
O modelo principal adotado é o Temporal Fusion Transformer (TFT), selecionado por sua adequação a cenários multivariados, previsão multihorizonte e mecanismos de interpretabilidade \cite{LIMETAL2021TFT}. A configuração experimental contempla modos de previsão pontual e quantílica, com priorização da previsão quantílica nos quantis \(p10\), \(p50\) e \(p90\) para análise de incerteza preditiva. Esse desenho permite avaliar simultaneamente tendência central e dispersão de risco, em linha com princípios de previsão probabilística \cite{GNEITING2007,KOENKERBASSETT1978}.

As configurações de treino incluem parâmetros de otimização, política de parada, definição de horizontes de avaliação e assinatura determinística de configuração para rastreabilidade. Cada execução registra identificadores determinísticos, metadados de ambiente e \texttt{parent\_sweep\_id} de coorte, de forma a permitir auditoria do experimento sem depender de memória operacional do pesquisador.

\section{Protocolo experimental (splits, seeds, folds)}
O protocolo experimental utiliza particionamento temporal explícito para treino, validação e teste, preservando a ordem cronológica das observações e evitando contaminação futura. Para reduzir sensibilidade a inicialização, as comparações entre configurações consideram múltiplas sementes e repetição de execuções dentro da mesma coorte experimental, identificada por \texttt{parent\_sweep\_id}.

A comparabilidade entre modelos segue regra de alinhamento temporal estrito por \textit{target timestamp} no momento de testes pareados. Esse critério evita comparações entre amostras não equivalentes e é pré-requisito para inferência estatística válida em testes de acurácia preditiva \cite{DIEBOLDMARIANO1995}. Os resultados confirmatórios deste trabalho serão obtidos sob pré-registro versionado, conforme política de \texttt{docs/08\_governance/} (\textit{preregistration\_round1}), e restritos a runs gerados após o reset operacional do silver/gold de 2026-05-10, conforme \texttt{docs/01\_architecture/ANALYTICS\_STORE\_ARCHITECTURE.md} (\textit{Archive pre-Phase B}).

\section{Baselines}
A comparação contra baselines explícitos é tratada como componente metodológico, não auxiliar. Os baselines admissíveis para este estudo são definidos canonicamente em \texttt{docs/04\_evaluation/BASELINES.md} e englobam: \textit{random walk}, AR(1), média histórica, EWMA-vol e quantis empíricos históricos. Esses estimadores são monotônicos por construção, o que assegura simetria de comparação em métricas probabilísticas (ver Seção 4.7).

A inclusão dos baselines no estudo segue duas regras de comparabilidade: (i) baselines são persistidos com o mesmo grão de \texttt{fact\_oos\_predictions} do modelo principal, permitindo análise pareada por \textit{target timestamp}; e (ii) baselines participam da mesma coorte (\texttt{parent\_sweep\_id}) e do mesmo \textit{scope} do candidato TFT \textit{all-features}, de modo que diferenças observadas reflitam modelagem e não diferenças de cobertura.

\section{Métricas de avaliação (pontuais e probabilísticas)}
As métricas pontuais adotadas incluem RMSE, MAE, acurácia direcional e viés médio, com apuração por horizonte e por split. Em paralelo, a avaliação probabilística considera \textit{pinball loss}, calibração marginal por quantil, PICP intervalar para o intervalo \(p10\)-\(p90\) e MPIW médio. Essa combinação busca evitar conclusões baseadas apenas em erro médio, incorporando qualidade da incerteza estimada \cite{GNEITING2007}.

A saída quantilica do TFT não garante, por construção, a monotonicidade \(p10 \leq p50 \leq p90\), patologia conhecida como \textit{quantile crossing} em regressão quantilica treinada independentemente. Para conciliar essa propriedade do estimador com a exigência de avaliação probabilística válida, o projeto adota a política canônica de variante quantilica documentada em \texttt{docs/04\_evaluation/METRICS\_DEFINITIONS.md} (categorização A/B/C) e detalhada em \texttt{docs/07\_reports/living-paper/20\_method.md} (\textit{Política de variante quantilica}): a pipeline aplica um \textit{guardrail} monotônico (rearranjo) antes da persistência, materializando colunas paralelas \textit{raw} e \textit{post-guardrail} no \textit{analytics store}. As hipóteses H1 (calibração) e H2a/H2b (\textit{pinball loss} contra baselines) são avaliadas sobre a variante \textit{post-guardrail} como primária, justificada pela validade do rearranjo monotônico sob perdas convexas simétricas em torno da mediana \cite{CHERNOZHUKOV2010REARRANGEMENT} e pelo enquadramento de regras de pontuação estritamente próprias \cite{GNEITING2007}. \textit{Checks} de patologia (intervalo negativo, degenerescência \(p10 = p90\)) permanecem sobre a variante \textit{raw} por construção semântica. Métricas de risco financeiro (\textit{Value-at-Risk}, \textit{Expected Shortfall}) são reportadas apenas \textit{post-guardrail}, dada a exigência de monotonicidade da função quantil em medidas de risco coerentes \cite{JORION2007,ACERBI2002}.

Para comparação entre modelos em ambiente fora da amostra, utiliza-se análise pareada de desempenho com teste de Diebold-Mariano \cite{DIEBOLDMARIANO1995} sob correção HAC e ajuste para amostras pequenas no estilo Harvey-Leybourne-Newbold \cite{HARVEY1997}, com ajuste de múltiplas comparações por Holm-Bonferroni, conforme \texttt{docs/04\_evaluation/STATISTICAL\_TESTS.md}.

\section{Explicabilidade e contribuição}
A caracterização da contribuição relativa de famílias de indicadores adota hierarquia metodológica explícita, documentada em \texttt{docs/04\_evaluation/EXPLAINABILITY.md}: \textit{Variable Selection Network} (VSN) nativa do TFT \cite{LIMETAL2021TFT} como evidência primária, importância por permutação como evidência complementar e ablação explicativa como verificação adicional. Conclusões de contribuição relativa exigem consistência em pelo menos dois dos três métodos; divergências entre métodos são interpretadas como instabilidade de contribuição preditiva, não como falha do estimador.

A análise é tratada como preditiva e condicional ao modelo treinado, não causal: VSN, permutação e ablação indicam o que o modelo usou, não o que move o retorno. Interpretações locais (por instância) são reportadas como ilustrativas; conclusões metodológicas são extraídas de evidência global por horizonte e da estabilidade entre sementes e \textit{folds}.

\section{Analytics Store e governança de experimento}
A governança experimental é operacionalizada por um repositório analítico com duas funções complementares: preservar dados primários de execução em camadas de persistência imutáveis e auditáveis, e materializar tabelas derivadas para análise comparativa e tomada de decisão. Nesse arranjo, identificadores de execução, assinaturas de configuração, \textit{fingerprints} de dataset/split, \texttt{parent\_sweep\_id} de coorte e status de execução são persistidos de forma estruturada para rastreabilidade integral do ciclo de modelagem.

A atualização do \textit{analytics store} ocorre sob escopo controlado por \textit{ScopeSpec}, conforme \texttt{docs/01\_architecture/ANALYTICS\_STORE\_ARCHITECTURE.md}, e a agregação de métricas derivadas adota \textit{groupby} consciente de coorte (\texttt{parent\_sweep\_id}), evitando que rankings, comparações pareadas e seleções de modelo misturem execuções de coortes heterogêneas. Como consequência metodológica, resultados de comparação podem ser reconstruídos a partir dos artefatos persistidos, sem necessidade de retreinar os modelos para cada rodada analítica. Essa propriedade fortalece auditabilidade, transparência e repetibilidade dos resultados reportados \cite{PINEAU2021REPRODUCIBILITY}.

\section{Critérios de qualidade e validação}
Os critérios de qualidade incluem validações de integridade temporal, completude mínima de cobertura, consistência de tipos e ausência de violações contratuais nas saídas analíticas. No caso de previsões quantílicas, aplica-se verificação de monotonicidade dos quantis (por exemplo, \(p10 \leq p50 \leq p90\)) sobre a variante \textit{raw}, além de checagens de unicidade das chaves temporais por execução e horizonte. Degenerescência quantilica (\(p10 = p90\)) acima de limiar declarado bloqueia a promoção do \textit{run} para análise confirmatória, conforme política operacional do projeto.

Para análise estatística pareada, somente são aceitas comparações com interseção temporal exata entre modelos, reduzindo risco de inferência inválida por coortes não comparáveis.

\subsection{Gates quantitativos de aceitação}
Para tornar auditáveis as decisões de promoção e descarte de \textit{runs}, o projeto operacionaliza os critérios de qualidade em quatro \textit{gates} explícitos com limiares quantitativos. Os limiares são contratos metodológicos declarados antes da observação dos resultados; sua versão canônica está em \texttt{docs/04\_evaluation/CALIBRATION\_AND\_RISK.md} e em \texttt{docs/07\_reports/living-paper/30\_results\_and\_analysis.md}.

O \textbf{Gate A} verifica o contrato quantilico bruto antes do guardrail. Aceita-se taxa de \textit{crossing} bruto menor ou igual a \(0{,}10\%\) por \textit{run}, taxa de \textit{crossing} pós-guardrail estritamente igual a zero (garantida por construção do rearranjo monotônico) e contagem de linhas com largura de intervalo negativa estritamente igual a zero. Falhas no Gate A indicam patologia do estimador na origem e bloqueam o uso do \textit{run} em comparação confirmatória, ainda que o pós-processamento mascare a violação a jusante.

O \textbf{Gate B} verifica o impacto do guardrail sobre as métricas probabilísticas, comparando variantes \textit{raw} e pós-guardrail por \textit{run}. Aceita-se \(\Delta\)pinball relativo menor ou igual a \(+1{,}0\%\), módulo do \(\Delta\)PICP menor ou igual a \(0{,}02\) e \(\Delta\)MPIW relativo menor ou igual a \(+5{,}0\%\). Esses limites garantem que a melhoria observada do candidato sobre os baselines não seja, em parte significativa, atribuível ao pós-processamento de rearranjo. Magnitudes acima dos limites limitam o claim ao estimador entregue (pipeline TFT mais guardrail), não ao TFT isolado.

O \textbf{Gate C} verifica robustez entre repetições da mesma configuração, medida pelo desvio-padrão da taxa de invalidação de quantis entre sementes e \textit{folds}. Um grupo de hiperparâmetros é classificado como instável quando o desvio-padrão de \texttt{invalid\_rate} excede \(0{,}005\); uma família de configurações é reprovada quando mais de \(20\%\) dos grupos forem instáveis. Esse critério penaliza candidatos cuja qualidade probabilística depende fortemente da semente, mesmo quando a média de erro é favorável.

O \textbf{Gate D} verifica sensibilidade a regime, exigindo que diferenças entre subgrupos de alta volatilidade e baixa volatilidade, ou entre regiões mais e menos fora da distribuição de treino, sejam reportadas com teste controlado a \(p < 0{,}05\) e sinal consistente entre indicadores complementares. Sua função é prevenir que conclusões médias agregadas mascarem degradação localizada da qualidade probabilística em janelas específicas.

A composição dos quatro \textit{gates} atua como barreira metodológica entre execução experimental e claim científico. Um \textit{run} é elegível para a análise confirmatória da Fase B quando aprovado nos Gates A e B; a promoção a evidência primária exige adicionalmente aprovação no Gate C; a leitura segmentada por regime do Capítulo 5 é qualificada pelo Gate D. Status \texttt{provisório} é atribuído após uma rodada de aprovação; status \texttt{oficial} requer aprovação em duas rodadas consecutivas, conforme política de governança do projeto.
